\chapter*{Özet}
\addcontentsline{toc}{chapter}{Özet}

%% Edit below this line
Bu projede, üretim ve montaj hatlarındaki video anomali tespiti (VAD), göreve özel hiçbir
eğitim gerektirmeden, Görü-Dil Modelleri (VLM) ile ele alınmıştır. İncelenen sahne, bir
parçanın konveyör üzerinde soldan sağa ilerleyişidir; amaç, kısa bir video klibinin
``normal'' mi yoksa bir ``anomali'' mi içerdiğine karar vermektir.

Üç farklı yöntem ailesi denenmiştir. İlki, modele doğrudan yapılandırılmış istemlerin
(prompt) verildiği temel yaklaşımdır; bu yaklaşım, modelin kliplere yersizce ``normal''
deme eğilimini açığa çıkarmıştır. İkincisi, yerelde çalıştırılan açık ağırlıklı bir model
(Qwen3-VL-8B) üzerine kurulan çok-persona oylama sistemidir: aynı klip, farklı bakış
açılarına sahip birden çok ``denetçi persona'' tarafından incelenir ve karar, çoğunluk
oyuyla verilir. Burada, modeli önce gördüğünü betimlemeye, sonra karar vermeye zorlayan
çıktı biçiminin, sistematik ``yanlış-normal'' hatalarını ortadan kaldırdığı görülmüştür.
Üçüncüsü, VERA yönteminden uyarlanan VERA-lite akışıdır; bu akış, modelin ağırlıklarını
değiştirmeden yönlendirici soruları iyileştiren bir öğrenen--eniyileyen döngüsünün yanı
sıra, etiket sızdırmayan (kör) bir değerlendirme yöntemi, farklı girdi temsilleri,
muhakeme bütçesi denetimi ve iki görüşü birleştiren bir topluluğu (ensemble) içerir.

Özel olarak hazırlanan $13$ kliplik kalem-konveyör veri setinde, sürekli video görüşü ile
ayrık yoğun-kare görüşünü birleştiren topluluk, iki görüşün birbirini tamamlayan kör
noktaları sayesinde $13/13$ doğruluğa ve sıfır yanlış pozitife ulaşmıştır. Tek başına
hiçbir temsil bu sonucu verememiştir; bu durum, zamansal (temporal) anomaliler için
sürekli video ile ayrık kare temsilleri arasında içsel bir ödünleşim bulunduğunu gösterir.
Koşuların tekrarlanması, tek bir görüşün gerçek doğruluğunun, şanslı bir tek koşuda
görülen $0.92$ yerine $0.846$ (üç koşunun da $11/13$'ü) olduğunu ortaya koymuştur. Ayrıca,
modelin muhakeme bütçesini sınırlamanın doğruluğu değiştirmeden gecikmeyi yaklaşık \%41
azalttığı ve klip bazındaki kararları tamamen kararlı (deterministik) hâle getirdiği
gözlenmiştir. Çalışmada olumsuz sonuçlar da dürüstçe raporlanmış; daha büyük veri setleri,
klasik bilgisayarla görü temel yöntemleri ve kare düzeyinde değerlendirme gibi
sınırlılıklar ile gelecek çalışmalar tartışılmıştır.

%% Until here
\vfill
%% Edit after {Anahtar Kelimeler:}
\textbf{Anahtar Kelimeler:} video anomali tespiti, görü-dil modelleri, sıfır-atış,
prompting, topluluk (ensemble), endüstriyel muayene.
\clearpage
